\section{Model}\label{sec:model}

\subsection{Workload and notation}\label{sec:notation}
A session has $T\sim\min(\mathrm{Geometric}(1/15),60)$ turns, so $E[T]=14.761$. Turn $k$ sends a prompt of $L_k=L_0+(k-1)(\Delta+d)$ tokens: a first prompt of $L_0=2000$ tokens, $\Delta=550$ tokens of tool output and $d=100$ decoded tokens per turn. Between turns the session waits for a tool for a time $Z$, independent across turns, exponential with mean 5\,s in the workload of record. A turn \emph{hits} if more than half of its prompt is served from cache; on the real engine a hit almost always reuses everything but the 550 new tokens and the last decoded token, so it prefills 551 tokens (all but 95 of the 185,613 hits in E2). First turns always miss, so in the long run the hit fraction is at most $1-1/E[T]=0.932$; finite windows can lie slightly above it (up to 0.94 in our runs). With every turn missing, a session's prefill work is $W(T)=2000\,T+325\,T(T-1)$ tokens with $E[W]=154{,}529$; with every follow-up hitting it is $2000+551(T-1)$, about 9,600 in expectation.

The engine has a KV pool of $M$ tokens. In the \emph{closed} system $n$ sessions are present and a finished session is replaced at once. The thrashing curve $\mu(n)$ is the long-run session completion rate (sessions/s); $\mu^*=\max_n\mu(n)$ is attained at $n^*$, and $\muinf=\lim_{n\to\infty}\mu(n)$ is the \emph{congested service rate}, a long-run rate of the closed system (an open system in congestion can complete sessions faster than $\muinf$ for a long time, \S\ref{sec:fluid}). The all-miss capacity $\lmiss$ is the throughput of the saturated engine with prefix caching disabled. In the \emph{open} system sessions arrive as a Poisson process of rate $\lambda$, and $N(t)$ counts sessions in the system (waiting, running or in a tool call). A session rate is estimated either as sessions completed per second or, with much lower variance because session work is heavy-tailed in $T$, as turns per second divided by $E[T]$ (``turns-based'').

\subsection{The thrashing curve for agent sessions}\label{sec:theory-curve}
Let $\Elb$ be the time-averaged KV footprint of a session over its lifetime, a length-biased mean in which long sessions and late turns weigh more (during the tool call after turn $k$ a session holds $L_k+d$ tokens). By the $H=\lambda G$ relation~\citep{brumelle1971relation,heyman1980relation}, the mean KV demand of $n$ present sessions is $n\Elb$. For our workload $\Elb$ is about 10.5k tokens when tool calls dominate the lifetime, so a 300k pool holds $\Nlb=M/\Elb\approx28.5$ sessions on average (weighting every turn equally, the last included, gives 10.6k and 28.4; we set the caps of \S\ref{sec:mech} from the latter). While $n\Elb$ is well below $M$, idle KV survives the tool call and $\mu(n)$ rises with $n$ through batching. Beyond it, admitting a waiting turn evicts idle KV, which belongs to sessions in tool calls or already waiting in the queue; they miss when their next turn is admitted. The peak $n^*$ lies below $\Nlb$ because resident demand fluctuates around its mean.

\subsection{The open system as a quasi-static approximation}\label{sec:fluid}
Replacing the departure rate of the open system by its closed-system value at the current population, the classical decomposition approximation (\S\ref{sec:background}), gives
\begin{equation}\label{eq:fluid}
\dot n=\lambda-\mu(n).
\end{equation}
Extend $\mu$ from integer $n$ to all $n\ge0$ by linear interpolation, so that it is continuous and piecewise linear with $\mu(0)=0$ and Eq.~\eqref{eq:fluid} has unique solutions, and assume that the stationary curve is unimodal: strictly increasing on $[0,n^*]$ to $\mu^*$, and strictly decreasing beyond $n^*$ with limit $\muinf<\mu^*$. For $0<\lambda<\mu^*$ let $n_1<n^*$ solve $\mu(n_1)=\lambda$. If $\lambda<\muinf$, $n_1$ is the only fixed point of Eq.~\eqref{eq:fluid}, and it attracts every trajectory. If $\muinf<\lambda<\mu^*$, there is one more, $n_2>n^*$, which is unstable, and from any $n>n_2$, $n\to\infty$: the good state is metastable. If $\lambda>\mu^*$ there is no good state. A solution from a congested state $x_0>n^*$ returns to $n_1$ if and only if $\mu>\lambda$ on $(n_1,x_0]$, the endpoint $n_1$ excluded because $\mu(n_1)=\lambda$; since $\mu>\lambda$ on $(n_1,n^*]$ and $\mu\ge\mu(x_0)$ on $[n^*,x_0]$, this holds if and only if $\lambda<\mu(x_0)$. As $x_0$ grows, $\mu(x_0)$ falls to $\muinf$, so $\muinf$ is the \emph{recovery threshold} from deep congestion. After a finite burst $x_0$ is finite, so recovery above $\muinf$ is possible; the threshold for a given burst depends on its size. Without the monotone tail, the recovery threshold from deep congestion is $\inf_{n\ge n^*}\mu(n)$ instead, and congested fixed points can exist (Appendix~\ref{app:proofs} proves these statements). Our measured curve is not monotone beyond $n^*$: in completed sessions it falls to 0.68 $\lref$ at $n=48$ and rises to 1.51 $\lref$ at $n=300$ (turns-based, 0.88 and 1.56 $\lref$; \S\ref{sec:curve}). Taken at face value, either estimate would put the recovery threshold from $n=300$ below $\lmiss$ and give a stable congested fixed point at 1.35 $\lref$. We attribute both features to the synchronized start, which no longer run has tested, and take $\muinf$ from Proposition~\ref{prop:A}.

In a finite system fluctuations can carry $N$ past $n_2$ without a burst, and an experiment sees only escapes within its horizon. We therefore use a measured \emph{from-empty edge} $\ledge$, the largest rate at which runs started from an empty engine stay stable over a stated horizon, and define the window
\begin{equation}\label{eq:window}
\mathcal W=\big(\muinf,\ \min(\mu^*,\ledge)\big).
\end{equation}
Inside $\mathcal W$ the good state exists from empty, and a burst that crosses the basin boundary leaves the system in a regime whose service rate is below the arrival rate. That regime is transient (the backlog grows), not a second stationary state, so we say metastable rather than bistable. In the terms of Bronson et al.~\citep{bronson2021metastable}, $\muinf$ is the hidden capacity, the load up to which the system heals itself after any trigger, and $\min(\mu^*,\ledge)$ the advertised capacity; at loads in $\mathcal W$ the good state is their vulnerable state, a burst is a trigger whose intensity matters, the congested regime is the metastable failure state, and the feedback of \S\ref{sec:background}, in which misses lengthen the queue and the gap between a session's turns, is its sustaining effect.

Eq.~\eqref{eq:fluid} assumes that the internal state (turn-index mix, cache contents, queue order) relaxes quickly compared with changes in $n$. An automated check of the model by an LLM agent, not independently reviewed (Appendix~\ref{app:prov}), found this reasonable at fixed points and for the sign of the drift in deep congestion, where the stationary $\mu$ is flat, but not near the cliff. The model is also not quantitative in congestion, and not only transiently: over the evaluation window the registered SGLang BURST runs still completed 1.13--1.35 $\lref$ at 300k and 1.44--2.03 times the native $\lref$, and their fitted backlog slopes were only 18--69\% and 68--78\% of the rate $\lambda-\lmiss$ that Eq.~\eqref{eq:fluid} gives with $\mu$ at its congested value $\lmiss$, taking $\lambda$ as the arrival rate each run received and $\lmiss$ as $\lref$, or 21--75\% and 70--80\% with $\lmiss$ 3.4\% above $\lref$ (\S\ref{sec:e2}, \S\ref{sec:e2native}, \S\ref{sec:props}). Under FCFS each waiting session advances one turn per cycle of the queue while new sessions keep arriving, so the backlog stays young and its turns are short (in the BURST windows served prompts were 9--20\% and 29--41\% below the stationary mean, \S\ref{sec:props}). At a token rate near the all-miss rate the engine then serves more turns per second than in the stationary mix, and since a session ends after each turn with probability 1/15 until its 60th, it also completes sessions faster than $\lmiss$, for as long as the backlog grows.

A work balance, which we do not make rigorous for the open system, suggests that the recovery threshold depends on the token rate in congestion, not on the departure rate. If hits vanish in congestion (as in every registered SGLang BURST run) and the engine computes prompt tokens at about $C=\lmiss E[W]$ (\S\ref{sec:props}) while they arrive at $\lambda E[W]$, then for $\lambda>\lmiss$ the uncomputed prompt work of the sessions in the system grows, and eventually so does their number, since each holds at most $W(60)$ tokens of it, while for $\lambda<\lmiss$ that work shrinks as long as the engine computes at $C$. If the zero-hit token rate is lower at some population that the system passes through on the way down, the threshold is correspondingly lower; in the fluid model it is $\inf_{n\ge n^*}\mu(n)$ rather than $\muinf$. On our engine the closed runs' zero-hit token rate was 4\% lower at $n=48$ than at $n=96$, with older turns (\S\ref{sec:props}). We therefore use Eq.~\eqref{eq:fluid} for the window, as a long-run statement for the recovery threshold, and for the capped system's drain time (\S\ref{sec:props}), not for burst sizes or backlog growth in congestion: in an earlier simulator with assumed costs (\S\ref{sec:sim}) the fluid critical burst $n_2-n_1$ at $2\lmiss$ was 20.5 sessions against an empirical 40, because burst sessions are fresh and small.

\subsection{Two results}\label{sec:props}

\begin{proposition}[Request-level FCFS, any online eviction rule]\label{prop:A}
Consider the closed system with $n$ sessions under any online eviction rule, one that does not foresee tool times, and assume that its long-run rates exist and are deterministic (Appendix~\ref{app:proofs}) and that:
\begin{enumerate}[label=(A\arabic*),leftmargin=*]
\item sessions share no cached tokens, each has at most one outstanding request, and evicted KV is discarded, not offloaded to another memory tier; tool times are i.i.d.\ with finite mean, independent of all other randomness, and unknown to the engine until they end;
\item admission is request-level FCFS: no request is admitted for the first time before a request that was already waiting when it joined the queue. More generally, it suffices that the order meets the condition of Appendix~\ref{app:proofs} for some $Q_0$ and $S_0$ independent of $n$, which allows reordering while at most $Q_0$ requests wait and lets a follow-up pass at most $S_0$ requests that were waiting when it joined. At most $B$ requests run at once, and at most $\alpha t+\beta$ admissions, re-admissions after a preemption included, occur in any interval of length $t$, with $B$, $\alpha$ and $\beta$ independent of $n$;
\item in congestion the engine computes prompt tokens no faster than with caching disabled: $\limsup_{n\to\infty}w_n\le C$, where $w_n$ is the long-run rate of prompt tokens that are computed rather than served from cache when their request is first admitted (each prompt token is counted once, so tokens recomputed after a preemption are not), and $C$ is the limit of that rate as $n\to\infty$ for the closed system with caching disabled.
\end{enumerate}
Then the long-run rates of cache-served prompt tokens and of hit turns tend to 0 as $n\to\infty$, so does the hit fraction if $\liminf_{n\to\infty}\mu(n)>0$, and $\limsup_{n\to\infty}\mu(n)\le\lmiss$. If also (A4) $w_n\to C$, that is, the congested engine stays saturated and computes prompt tokens as fast as with caching disabled, then $\mu(n)\to\lmiss$, so the hit fraction tends to 0.
\end{proposition}
\begin{proof}[Proof sketch]
By (A1), the $r_i$ prompt tokens that follow-up $i$ reuses at its first admission are its session's own and stay in the pool from the end of the session's previous turn until that admission, an interval no shorter than its queue wait $D_i$, and the token-time that these intervals claim does not overlap. Resident token-time accrues at rate at most $M$, so follow-ups that wait at least $K$ reuse tokens at rate at most $M/K$, whatever the eviction rule. By (A2) a follow-up that waits less than $K$ joined behind fewer than $\alpha K+\beta+Q_0+S_0$ requests, so just before it joined more than $n-B-\alpha K-\beta-Q_0-S_0$ sessions were in tool calls; Lemma~\ref{lem:waits} (Appendix~\ref{app:proofs}), which uses that the engine does not see tool times, shows that such follow-ups occur at a rate of order $1/n$, and each reuses at most $L_{\max}=40{,}350$ tokens. Letting $n\to\infty$ and then $K\to\infty$, the cache-served token rate tends to 0, and since a hit reuses more than $L_0/2$ tokens, so does the rate of hits. Counted once, at its request's first admission, each of a session's $W$ prompt tokens is either served from cache or computed, so $w_n$ is $\mu(n)E[W]$ less the cache-served token rate, and with caching disabled the same count gives $C=\lmiss E[W]$; (A3) gives the bound and (A4) the limit. The same count gives turns at the rate $\mu(n)E[T]$, so the hit fraction tends to 0 whenever $\mu(n)$ stays bounded away from 0.
\end{proof}

The proof's last step, from vanishing reuse to $\lmiss$, is the bound that MISA-T notes at prefill saturation, a request rate of at most the prefill throughput over the mean uncached prompt tokens per request~\citep{hong2026misat}; what the proposition adds is that request-level FCFS makes reuse vanish in congestion whatever the online eviction rule. (A2) describes an admission pass that stops at the first request that does not fit, as in our simulators and, on our reading of their sources, in SGLang and vLLM under their default FCFS policies (Appendix~\ref{app:proofs}); vLLM also preempts running requests and returns them to the head of its queue, which (A2) allows, because its ordering clause concerns first admissions; with asynchronous scheduling it also lets a follow-up pass a bounded number of just-preempted requests, so it meets the general condition of (A2) with $Q_0=0$ and a bounded $S_0$ (Appendix~\ref{app:proofs}). A diverging mean cycle time would not be enough: Lemma~\ref{lem:waits} uses bounded admissions and i.i.d.\ tool times to bound the tail of the number of sessions in tool calls independently of $n$, which makes it rare that almost all sessions are in tool calls when one returns. The lemma's bound is of order $1/n$ but has large constants, so the proposition describes the limit, not the populations we ran. The lemma needs FCFS order only while more than $Q_0$ requests wait, so, on our reading of its source, the proposition also covers SGLang's stock LPM, which stops re-sorting its queue whenever more than 128 requests wait, so that new requests queue behind all waiting ones ($Q_0=128$). It does not cover program-level FCFS, which puts an old session's follow-up ahead of newer sessions' requests however long the queue: in our earlier simulator its hit fraction stays near 0.03 from $n=64$ to 300, while FCFS has none. A pure resident-first order, such as LPM without the fallback, violates (A2) outright: a returning session whose KV is still resident is admitted ahead of the queue.

(A1) holds for our random-token sessions, apart from rare one-token coincidences (5 of the 19,924 first turns in E2 had one cached token), but not for agents that share a system prompt: one resident copy then serves every session. We expect the argument to apply to the private suffixes, with $\lmiss$ and $C$ measured with the shared prefix cached, but have not shown this. (A1) also rules out offloading evicted KV to host memory, which would let a follow-up reuse tokens that held no pool space while it waited. (A3) and (A4) treat the engine as a server of prompt tokens, as E0a suggests (\S\ref{sec:edge}). They concern long-run rates of the closed system, in which every order serves turns in the stationary mix, since each completed session contributes all its turns; our measurement windows did not have that mix, and the token rate varied between them, so we can check (A3) and (A4) only approximately. We estimate $C$ by E0a's uncached token rate at 32 sessions, 11,243 tokens/s. With zero hits, closed runs computed prompt tokens at 1.03~$C$ at $n=96$, whose served prompts were within 3\% of the stationary mean, on SGLang and on vLLM (against vLLM's own $C$), at 0.99~$C$ at $n=48$ (prompts 12--14\% longer) and at 1.10--1.11~$C$ at $n=300$ (28--30\% shorter); the registered SGLang BURST runs computed 1.01--1.03~$C$ at 300k (prompts 9--20\% shorter) and 1.07--1.12 times the 32-session all-miss rate of the native pool (prompts 29--41\% shorter). So on our engine (A3) and (A4) hold only approximately: to within a few percent at the stationary mix, and to within 12\% over the windows we measured. We estimate $\lmiss$ by $\lref$, E0a's turns-based rate at 32 sessions (0.0715\sps); $\lref$ and $C/E[W]$ estimate the same cache-disabled limit. $C/E[W]$ is 0.0728\sps, 1.7\% above $\lref$, because the turns in E0a's window were longer than the stationary mix (10,648 against 10,469 prompt tokens on average). E0a's 64-session check computed 1.7\% more tokens per second than the 32-session runs and 3.2\% more than the same seed at 32 sessions, although its window held longer turns (11,456 prompt tokens on average), so the 32-session reference falls short of saturation by about 2--3\%, perhaps more. With the 64-session rate in place of $C$, $C/E[W]$ is 1.034 $\lref$, and with the zero-hit rate at $n=96$ (11,602 tokens/s) it is 1.050 $\lref$. Under (A4) the proposition therefore puts $\muinf$ about 2--5\% above $\lref$, an offset within the standard error of $\lref$ over seeds (0.0038\sps, 5\%); on vLLM the three estimates are 1.020, 1.038 and 1.050 $\lref$ ($\lref=0.0657$\sps).

The bound holds for every online eviction rule, including protection of queued prefixes and TTL pins that keep request-level FCFS order, where (A4) can fail. In our earlier simulator, under FCFS order, a hard pin of queued prefixes gives 0.919 $\lmiss$ at $n=300$ (\S\ref{sec:mech}) but 1.04--1.15 $\lmiss$ at $n=48$--150 (3 seeds), as a limit statement allows, and TTL pins of 5--20\,s, which that simulator keeps past expiry while the session's next request waits, give 0.10--0.29 $\lmiss$ at $n=300$ with the earlier release rule of \S\ref{sec:mech}, with the GPU busy only 8--33\% of the time (3 seeds). Both still hit 13\% and 10--23\% of turns at $n=300$, so a finite-$n$ hit fraction does not by itself separate covered orders from uncovered ones. The proposition makes their hit rate vanish as $n$ grows, and their hit fraction too unless their throughput tends to 0, as the TTL pins' may; the hard pin keeps the GPU busy, and its hit fraction falls from 0.23 at $n=48$ to 0.13 at $n=300$. The TTL rows of \S\ref{sec:mech} that exceed $\lmiss$ also serve pinned programs first, an order that (A2) excludes. Belady's offline rule~\citep{belady1966study} foresees tool times and is outside (A1): the token-time step of the proof covers it, Lemma~\ref{lem:waits} does not.

\begin{remark}[resident-first orders; a conjecture]
Under pure LPM (without SGLang's 128-request fallback) in congestion we expect the follow-up hit fraction to tend to $P(Z<R_\infty)$, with $R_\infty$ the congested lifetime of idle KV, so $\muinf^{\mathrm{LPM}}$ exceeds $\lmiss$ through tool calls shorter than $R_\infty$; the simulator agrees (\S\ref{sec:mech}), but we have no proof.
\end{remark}

\begin{proposition}[cap identity; elementary]\label{prop:cap}
Let a FIFO gate admit at most $c$ sessions, each holding its slot from admission to the end of its last turn, tool calls included. In the closed system with $n$ sessions the admitted sessions form a closed system of $\min(n,c)$ sessions, so $\mu_{\mathrm{cap}}(n)=\mu(\min(n,c))$.
\end{proposition}
\begin{proof}
When an admitted session finishes, the gate admits its head at once and the replacement joins the gate.
\end{proof}

With $c\le n^*$, $\mu_{\mathrm{cap}}$ is non-decreasing, so under Eq.~\eqref{eq:fluid} the capped system has one attracting state if and only if $\lambda<\mu(c)$. A burst of $b$ sessions that arrives when $n_1<c$ sessions are present fills the $c-n_1$ free slots at once, so the gate empties after about $(n_1+b-c)/(\mu(c)-\lambda)$. The registrations' feasibility check used the cruder $b/(\mu(c)-\lambda)$, which ignores the free slots: 1571\,s at 1.35 and 1889\,s at 1.65 $\lref$ for the registered test, and 668\,s and 1043\,s at 2.5 and 4.0 $\lref$ for the native one. With the populations present at the burst (5--12 sessions at 300k, 11--32 with the native pool), the fluid estimate is 1453--1508\,s and 1757--1804\,s, and 491--518\,s and 840--876\,s (compared with the engine in \S\ref{sec:e2} and \S\ref{sec:e2native}). This is the standard closed-system equivalence of a system with a bounded multiprogramming level (\S\ref{sec:background}). The stochastic counterpart (positive recurrence for $\lambda<\mu(c)$) should follow from standard arguments; we have not written them out. A cap $c>n^*$ raises the congested level to $\mu(c)$ but leaves the window $(\mu(c),\mu^*)$ open.
